\documentclass[11pt]{article}
\usepackage[margin=1.15in]{geometry}
\usepackage{amsmath,amssymb}
\usepackage{booktabs}
\usepackage{graphicx}
\usepackage{longtable}
\usepackage[table]{xcolor}
\usepackage{caption}
\usepackage[colorlinks=true,linkcolor=blue,citecolor=blue,urlcolor=blue]{hyperref}
\usepackage{titlesec}
\usepackage{enumitem}
\usepackage{fancyhdr}

\captionsetup{font=small,labelfont=bf}
\titleformat{\section}{\normalfont\large\bfseries}{\thesection}{0.6em}{}
\titleformat{\subsection}{\normalfont\normalsize\bfseries}{\thesubsection}{0.6em}{}
\setlist{itemsep=2pt,parsep=0pt,topsep=3pt}

\pagestyle{fancy}
\fancyhf{}
\rhead{\small Kahlus / Neural-CASP --- status, 25 July 2026}
\cfoot{\thepage}
\renewcommand{\headrulewidth}{0.4pt}

\newcommand{\bits}{\,\mathrm{bits}}

\begin{document}

\begin{center}
{\LARGE\bfseries A leakage-controlled test for incremental\\[2pt]
predictive information in neurophysiological signals}\\[10pt]
{\large Current state, validated apparatus, and complete findings inventory}\\[16pt]

{\normalsize Kahlus / Neural-CASP}\\[6pt]
{\small Prepared 25 July 2026 for Todd P.\ Coleman and Amrith Lolitkar}\\[4pt]
{\small Repository \texttt{Kahlus-V1}, branch \texttt{fix/forecast-overlap-guard}}
\end{center}

\vspace{4pt}
\hrule
\vspace{10pt}

\begin{quote}
\small
\textbf{Statement of scope.} This document reports what the Kahlus evaluation
apparatus has been demonstrated to do, with every statistic traceable to a named
artifact in the repository. \textbf{It reports no positive biological finding,
because none exists.} Every real-data result obtained to date is a negative or an
underpowered null, and each is labelled as such below. The results we consider
substantive are methodological: a calibrated detection floor
(Section~\ref{sec:calib}), a quantified evaluation-leakage failure mode
(Section~\ref{sec:leak}), and one instance in which the control battery rejected a
result that the bootstrap confidence interval alone would have certified
(Section~\ref{sec:caught}). We regard the last of these as the most useful single
result the project has produced.
\end{quote}

\vspace{6pt}
\hrule
\vspace{12pt}

\section{The estimand}

For an outcome $Y$ at horizon $h$, a nuisance block $B$, and a candidate block
$Z$, we define the \emph{residual forecastability score}
\begin{equation}
\mathrm{RFS}_{\text{bits}}
= \frac{\mathrm{NLL}(Y \mid B) - \mathrm{NLL}(Y \mid B, Z)}{\ln 2},
\label{eq:rfs}
\end{equation}
where $\mathrm{NLL}$ is the mean out-of-fold binary cross-entropy of a calibrated
probability forecast, with probabilities clipped to $[10^{-5}, 1-10^{-5}]$.
Implementation: \texttt{forecastability/m1.py::\_rfs\_bits}.

Equation~\eqref{eq:rfs} is the out-of-sample log-likelihood gain from adding $Z$
to $B$, in bits per anchor. Methodologically it is a cross-fitted,
subject-held-out analogue of nested $\Delta R^2$ / $\Delta \mathrm{CV}R^2$ model
comparison --- the same incremental-value question posed in
Rao et al.~(2025, \S2.5.2), evaluated out-of-sample under a mandatory control
battery rather than in-sample.

\paragraph{What RFS is not.} RFS is \emph{not} a mutual-information estimate, and
we do not interpret it as one. McAllester and Stratos (2018) prove that no
distribution-free high-confidence lower bound on mutual information from $N$
samples can exceed $O(\log N)$; any estimator read as a certified MI bound
inherits that limitation. RFS is a difference of out-of-fold mean log-losses
between two nested models with bounded predictions --- a bounded functional of
held-out forecasts. Concentration and conformal arguments apply to \emph{that}
quantity. Reading RFS as a bound on $I(Y; Z \mid B)$ is not licensed by our
procedure.

\medskip
The blocks are fixed across all findings below:
\begin{itemize}
  \item $B$ (nuisance) always contains the event base rate, recent event history,
        and cortical spectral features. A positive RFS therefore means $Z$ carries
        information the cortex's own recent state does not already supply.
  \item $Z$ (candidate) is the block under test.
  \item $Y$ is a future event indicator at horizon $h$.
\end{itemize}

\section{Evaluation protocol}

\subsection{Cross-fitting and clustering}
Nuisance and residual models are fit out-of-fold with \texttt{GroupKFold} at the
\emph{subject} level. All resampling is at the cluster (subject) level, never
per-anchor, since anchors within a recording are strongly dependent.

\subsection{Baseline ladder}
The residual is scored against the \emph{minimum-NLL} member of a candidate
baseline set (moving average, persistence, base rate, cross-fitted nuisance-only
model). Reporting a gain over a weak baseline is a standard route to inflated
incremental-value claims; selecting the strongest baseline removes that degree of
freedom.

\subsection{Control battery}
Every candidate result must survive all of:

\begin{center}
\small
\begin{tabular}{@{}ll@{}}
\toprule
Control & What it breaks \\
\midrule
Label shuffle & Any association at all \\
Time shift & Temporal alignment \\
Within-subject circular shift of $Z$ & Alignment, preserving $Z$'s autocorrelation \\
Subject-identity probe & $Z$ acting as a subject fingerprint \\
Cluster permutation & Supplies the $p$-value \\
\bottomrule
\end{tabular}
\end{center}

\subsection{Certification rule}
\label{sec:cert}
Implemented in \texttt{\_rfs\_eval.py::horizon\_real\_passes}. An edge is
certified only if \emph{all} of the following hold:
\begin{align*}
\mathrm{RFS}_{\text{bits}} &> 0.02 \\
\mathrm{LCB} &> 0 \quad \text{(one-sided 95\% subject-cluster bootstrap)} \\
\text{every control} &< 0.40 \times \mathrm{RFS}_{\text{bits}} \\
\text{nuisance probe failures} &= \varnothing
\end{align*}
plus minimum subject and minimum positive-event counts. The $0.02\bits$ floor is a
detection threshold, not a claim of clinical utility.

\section{Validated apparatus: calibration}
\label{sec:calib}

Detection power is a rate over seeds; a single run per condition cannot estimate
it. We inject the candidate coupling at graded strengths into a synthetic fixture
with known ground truth and run the complete gate at each strength across
24 independent seeds.

\begin{table}[h]
\centering
\small
\caption{Detection-power sweep, 24 seeds per point, synthetic ground truth.
Detection $=$ fraction of seeds whose one-sided 95\% bootstrap lower bound clears
zero. Source: \texttt{runs/autonomic\_calibration/detection\_power\_sweep.csv}.}
\begin{tabular}{@{}lccc@{}}
\toprule
Injected coupling & Detection power & Detections & Median RFS (bits) \\
\midrule
0.0 (null anchor) & \textbf{0.042} & 1 / 24 & $-0.0005$ \\
0.5 & 0.583 & 14 / 24 & $+0.0084$ \\
1.0 & \textbf{0.792} & 19 / 24 & $+0.0211$ \\
2.0 & 0.667 & 16 / 24 & $+0.0145$ \\
\bottomrule
\end{tabular}
\end{table}

\paragraph{False-positive rate.} The null anchor fires on 1 of 24 seeds,
$= 4.2\%$, against a nominal $5\%$ for a one-sided 95\% lower bound. The single
false positive (seed 2) has $\mathrm{RFS} = 0.0022\bits$ with
$\mathrm{LCB} = 0.00023$ --- barely clearing zero, which is the expected
appearance of a correctly calibrated borderline detection rather than a
malfunction. Per-seed values are retained in
\texttt{detection\_power\_per\_draw.csv} so individual detections can be audited.

\paragraph{Non-monotonicity is real and is explained.} Power declines from 0.792
at strength 1.0 to 0.667 at strength 2.0, and median bits from $+0.0211$ to
$+0.0145$. This persists across all 24 seeds in both panels and is not seed noise.
RFS is a \emph{residual}: once the outcome becomes highly predictable from base
rate and recent history alone, $B$ absorbs the signal and the residual
attributable to $Z$ shrinks. Detection power is therefore the calibration
quantity we report; median recovered bits is shown alongside it and is explicitly
not a calibration curve.

\paragraph{Scope.} These figures are synthetic by construction. They establish
that the estimator is calibrated before it is pointed at real data. They are not
evidence about any cohort.

\section{Validated apparatus: a quantified leakage failure mode}
\label{sec:leak}

The project's original headline result was an artifact, and the mechanism is
worth reporting because it survives an otherwise correct subject-held-out split.

The historical forecasting task took a 128-sample window, used
$X = \mathrm{seq}[0{:}127]$ as input and $Y = \mathrm{seq}[1{:}128]$ as target,
and scored mean squared error over the \emph{full} 127-length target. Input and
target therefore shared \textbf{126 of 127 samples}. A model that copies its input
shifted by one is exact on 126 positions; the single genuinely future sample is
averaged away.

\begin{table}[h]
\centering
\small
\caption{Identical ridge model, identical real Sleep-EDF data, two scoring rules.
Source: \texttt{docs/paper/ieee\_figure\_source/data/amrith\_overlap\_headline.json}.}
\begin{tabular}{@{}lc@{}}
\toprule
Scoring rule & Ridge MSE \\
\midrule
Full 127-sample target (126/127 shared with input) & $0.00180$ \\
Isolated strictly-future sample & $0.22694$ \\
\midrule
\textbf{Inflation factor} & $\mathbf{126\times}$ \\
\bottomrule
\end{tabular}
\end{table}

\paragraph{Distinction from the existing leakage literature.} Brookshire et
al.~(2024) and Del Pup et al.~(2025) document \emph{train/test} dependence:
related windows crossing folds. Reported inflation in that literature is of order
8--13\% F1 (Dehghani et al., 2019) up to 28\% accuracy (Kattenborn et al., 2022).
The failure above is different in kind --- it is \emph{within-window} input/target
overlap, which survives a perfectly correct subject-held-out split and is
invisible to any split-level audit. That is why its magnitude is two orders larger.

\paragraph{Remediation, and how it was verified.} Forecasting tasks now carry a
boolean \texttt{metric\_mask} true only on target positions absent from the input
window, applied to both the training loss and the evaluation metric. The
regression test constructs a predictor that exploits the overlap: it scores
$0.059$ unmasked and $0.470$ masked, an \textbf{8$\times$ discriminating gap}. It
is that gap, not agreement between two metrics, that demonstrates the mask
functions.

A guard (\texttt{data/forecast\_contract.py}) makes the unsafe configuration
unconstructible: a forecasting task is valid only if
$\text{horizon} \geq \text{window length}$ or a non-empty strictly-future mask is
attached. The historical claim is marked invalid in place, not deleted.

\paragraph{Person-level split auditing.} A separate audit addresses a distinct
failure: declared-\texttt{subject\_id} checks trust their labels, so when two
nights of one person are labelled as two subjects the check passes while the
person sits on both sides of the split. The regression test constructs exactly
that manifest and asserts that the conventional check \emph{passes} while the
person-level check \emph{fails}. Where person identity cannot be established the
audit fails closed.

\section{Result: the control battery rejects a bootstrap-certifiable effect}
\label{sec:caught}

This is, in our assessment, the most substantive result the project has produced.

We ran the gate on MIT-BIH Polysomnographic Database (slpdb) --- 11 records,
9 subjects, 6152 windows, claim-grade bootstrap, subject-held-out --- with
$Z$ an autonomic block (QRS-derived heart rate, blood-pressure and respiration
log-power envelopes), $B$ cortical spectral features plus event history, and
$Y$ a future arousal indicator. Arousals were stratified into respiratory
(hypopnea/apnea-terminating) and non-respiratory (bare and leg-movement) classes,
because the precedence literature diverges sharply between them.

\begin{table}[h]
\centering
\small
\caption{slpdb autonomic gate, by arousal class and horizon. ``Worst control''
is the largest of the time-shift and circular-shift surrogates, expressed as a
multiple of the real effect; the certification rule requires $< 0.40$.}
\begin{tabular}{@{}llrrrcc@{}}
\toprule
Class & $h$ & Pos.\ events & RFS (bits) & LCB & Worst ctrl / real & Subj.\ probe \\
\midrule
Non-respiratory & 1 & 190 & $-0.0265$ & $-0.1186$ & --- & fail \\
Non-respiratory & 2 & 288 & $-0.0139$ & $-0.0519$ & --- & fail \\
Non-respiratory & 4 & 408 & $-0.0236$ & $-0.0502$ & --- & fail \\
\midrule
Respiratory & 1 & 1360 & $+0.0058$ & $-0.0022$ & $1.42$ & fail \\
Respiratory & 2 & 1795 & $+0.0087$ & $\mathbf{+0.0008}$ & $\mathbf{1.05}$ & fail \\
Respiratory & 4 & 2170 & $+0.0111$ & $\mathbf{+0.0015}$ & $\mathbf{1.57}$ & fail \\
\bottomrule
\end{tabular}
\end{table}

\paragraph{The respiratory rows would have been reported as positive.} At $h=2$
and $h=4$ the subject-cluster bootstrap lower bound clears zero, and the cluster
permutation test returns $p = 0.031$ at $h=2$ (512 permutations, patient-cluster
sign-flip, exact mode). On the confidence interval and $p$-value alone, these are
findings.

\paragraph{They are rejected, on two independent grounds.}
\begin{enumerate}
  \item \textbf{The time-shift control scores at or above the real effect at every
  horizon} --- $1.42\times$, $1.05\times$, $1.57\times$, against a required
  $<0.40\times$. Destroying temporal alignment does not reduce the score. Whatever
  is being measured is therefore not predictive in time.
  \item \textbf{The subject-identity probe fails at every horizon in both
  classes.} The autonomic block identifies which subject a window came from. With
  9 subjects of differing apnea severity and correspondingly different arousal
  base rates, a subject fingerprint alone manufactures an apparent residual.
\end{enumerate}

The effect is static, subject-level structure --- not forecasting. \textbf{The
bootstrap confidence interval would have certified it; the control battery
rejected it.} We report this as direct real-data evidence that the controls are
load-bearing rather than ceremonial, and as a concrete instance of what a
CI-only pipeline would have shipped as a finding.

\paragraph{The non-respiratory arm is underpowered, not null.} That class has 190
positive events at $h=1$, and \textbf{73\% of all non-respiratory arousals in the
cohort (139 of 191) come from a single subject}; only 4 of 9 subjects have
$\geq 10$ such events. The $h=1$ interval spans $[-0.119, +0.043]$ --- consistent
with anything from a large negative to a moderate positive. This is the class for
which autonomic precedence is actually supported in the literature, and slpdb
cannot test it. We report it as underpowered and draw no inference.

\section{Complete findings inventory}

Every real-data finding, with its status stated exactly.

\begin{table}[h]
\centering
\footnotesize
\caption{All findings. No entry is a positive biological result.}
\begin{tabular}{@{}llp{6.6cm}@{}}
\toprule
ID & Status & Result \\
\midrule
F1 & Invalidated artifact &
Overlap-inflated forecasting MSE; $126\times$, marked invalid in place
(Section~\ref{sec:leak}). \\
\addlinespace
F2 & Negative (architecture) &
Isolated strictly-future forecasting; learned models separate clearly from
trivial baselines, but the GRU does not meaningfully beat a 127-tap ridge and
loses at $h=64$. See Table~\ref{tab:f2}. \\
\addlinespace
F3 & Scout-grade negative &
Peripheral channels $\to$ future stage transition, Sleep-EDF. Does not survive
controls at scout thresholds. Design-limited: Sleep-EDF cassette carries no clean
ECG. \\
\addlinespace
F4 & \textbf{Powered negative} &
Complexity beyond spectral, 78 subjects, 413{,}828 windows, $n_{\text{boot}}=2000$.
Wake $-0.330$ $[-0.349,-0.313]$; NREM $-0.284$ $[-0.297,-0.271]$;
REM $-0.161$ $[-0.177,-0.146]$. All CIs exclude zero on the negative side: the
complexity block actively hurts. \\
\addlinespace
F5 & \textbf{Powered negative} &
Propofol sedation, OpenNeuro ds005620, full 21/21 cohort, 2527 windows.
RFS $-0.2072$ $[-0.3084,-0.1208]$. Same conclusion as F4 under a different
etiology. Not TMS-evoked PCI; no consciousness or clinical claim. \\
\addlinespace
F6 & Built, unrun &
NSRR MESA/SHHS. Synthetic known $+0.021$ / null $-0.000$ passes; real cohort
blocked on credentialed access. \\
\addlinespace
F7 & Negative $+$ underpowered &
slpdb, Section~\ref{sec:caught}. Respiratory: caught false positive.
Non-respiratory: underpowered. \\
\bottomrule
\end{tabular}
\end{table}

\begin{table}[h]
\centering
\small
\caption{F2 --- isolated single-future-sample forecasting MSE, subject-held-out
Sleep-EDF, every model trained directly per horizon. Lower is better; best per
row in bold. Source: \texttt{runs/amrith\_isolated\_check/sweep.csv}.}
\label{tab:f2}
\begin{tabular}{@{}lccccc@{}}
\toprule
$h$ (samples @100\,Hz) & Random & Mean-of-trace & Persistence & Ridge-AR & GRU \\
\midrule
1  & 1.855 & 1.037 & 0.281 & 0.227 & \textbf{0.211} \\
2  & 1.865 & 1.081 & 0.338 & 0.310 & \textbf{0.277} \\
4  & 1.981 & 1.155 & 0.488 & 0.430 & \textbf{0.400} \\
8  & 2.070 & 1.232 & 0.853 & 0.675 & \textbf{0.583} \\
16 & 2.054 & 1.264 & 1.334 & 0.822 & \textbf{0.690} \\
32 & 2.254 & 1.432 & 1.855 & 0.956 & \textbf{0.883} \\
64 & 2.208 & 1.395 & 2.084 & \textbf{0.884} & 0.913 \\
\bottomrule
\end{tabular}
\end{table}

\paragraph{Reading F2 correctly.} Both learned models separate decisively from
random and mean-of-trace at every horizon, and the margin over persistence widens
with $h$ --- at $h=64$ persistence has degraded to 2.084 while ridge holds 0.884.
That is a real effect. But the GRU's advantage over a regularised linear
autoregression is small and reverses at $h=64$. At single-channel millisecond
scale the signal is close to linear, and a ridge captures nearly all of it. This
is consistent with the finding that well-regularised linear and Riemannian
methods remain competitive on EEG (Lotte et al., 2018), and with recent
foundation-model evaluations reporting no consistent advantage for larger
architectures in low-data or transfer regimes.

\section{What the negatives do and do not establish}

We state the limitations plainly, because several materially weaken the
interpretation of the results above.

\begin{enumerate}
  \item \textbf{No positive real-data finding exists.} F2--F5 are negative; F6 is
  unrun; F7 is negative in one class and underpowered in the other.
  \item \textbf{The powered negatives are on targets where a null was the expected
  outcome.} Sleep-stage discrimination is well served by spectral features, so
  complexity adding nothing is unsurprising. This weakens the rhetorical force of
  F4 and F5 even though it does not weaken their validity.
  \item \textbf{F4 and F5 test one specific complexity block}, not ``complexity''
  in general.
  \item \textbf{There is a published null against the broad autonomic
  hypothesis.} Rao et al.~(2025) tested cardiac autonomic predictors in a nested
  design and found $\Delta R^2 < 0.03$, $p > 0.25$ (HRV RMSSD partial
  $r = 0.16$, $p = 0.29$) for subjective sleep quality. Our defence is timescale
  and outcome specificity --- a within-night, event-locked arousal at a fixed
  horizon rather than a night-level subjective rating --- and that is a defence,
  not a refutation.
  \item \textbf{slpdb is an apnea cohort}, so for the respiratory class it is the
  sub-population where autonomic precedence is \emph{least} supported. It is an
  adversarial case, not a small version of MESA.
  \item \textbf{Calibration evidence is synthetic.} It bounds the estimator's
  operating characteristic under a generative model we wrote. It is not external
  validation.
  \item \textbf{RFS is not monotone in effect size} (Section~\ref{sec:calib}), so
  detection power, not recovered bits, is the calibration quantity.
\end{enumerate}

\section{On the relationship to Rao et al.\ (2025)}

We wish to be precise about this, since we cite the work as methodological
grounding and do not wish to misrepresent it.

The positive result in Rao et al.\ is \emph{intrinsic} infraslow rhythmicity ---
EGG ISO amplitude variance predicting subjective sleep quality
($\Delta R^2 = 0.13$, $p = 0.014$, $n = 49$ nights). The \emph{coupling} measures
added essentially nothing ($\kappa$: $\Delta R^2 = 0.0005$, $p = 0.84$; parietal
MI: $\Delta R^2 = 0.046$, $p = 0.17$), and the paper's own conclusion attributes
the effect to consistency of infraslow modulation within each organ rather than
to stomach--brain coupling. Citing that paper as motivation for a \emph{coupling}
claim would misread it, and we do not.

What we take from it is the nested-model design of \S2.5.2: a baseline model of
macro-architectural and covariate features, then the candidate block added and
partial $R^2$ reported. Our difference is cross-fitting and subject-held-out
evaluation in place of in-sample fit, plus a mandatory control battery. We note
that the paper already reports a cross-validated $\Delta \mathrm{CV}R^2$
($-0.015$ for the phase-coupling measure), so the CV-shrinkage concern is
acknowledged in that work.

\section{Immediate requirements}

\begin{enumerate}
  \item \textbf{NSRR credentialed access (MESA, then SHHS).} slpdb has settled
  nothing about the F6/F7 question, for the concrete reason given in
  Section~\ref{sec:caught}: the arousal class the hypothesis concerns has 190
  events, 73\% from one subject. MESA is required, not preferable. SHHS would then
  provide a dataset-held-out replication --- an axis almost nobody evaluates, and
  where the published cross-cohort drop is 0.12--0.32 absolute F1
  (Ehrlich et al., 2024).
  \item \textbf{Separation from arousal \emph{detection}.} ECG-only arousal
  detection already reaches AUROC 0.93 on MESA (Li et al., 2020) and AUC 0.91 from
  five post-onset beats (Basner et al., 2007). Our estimand is strictly-future
  forecasting conditioned on a cortical baseline. Any presentation that does not
  make this explicit will be read as a weaker version of existing detection work.
\end{enumerate}

\section{Reproduction}

\small
\begin{verbatim}
python3.11 -m venv .venv && .venv/bin/pip install -e '.[figures,moabb]' pytest edfio

PYTHONPATH=src python -m pytest -q            # 893 passed, 4 skipped, 114 subtests
uvx ruff@0.16.0 check .                       # clean

# calibration (Section 3)
PYTHONPATH=src python3 scripts/autonomic_rfs_calibration_sweep.py \
    --out runs/autonomic_calibration --seeds 24 --strengths 0 0.5 1 2

# slpdb gate (Section 5)
PYTHONPATH=src python3 scripts/run_slpdb_autonomic_gate.py \
    --out artifacts/slpdb_autonomic_non_respiratory --bootstrap-mode claim \
    --arousal-class non_respiratory --data-dir <local slpdb>
\end{verbatim}
\normalsize

Every gate emits \texttt{claim\_scope} and \texttt{stop\_reason} into its JSON and
Markdown artifacts, so a run's permitted interpretation travels with its numbers.

\section*{References}
\small
\begin{itemize}[leftmargin=1.2em,label={}]
\item Basner, M., et al. (2007). An ECG-based algorithm for the automatic
identification of autonomic activations associated with cortical arousal.
\emph{Sleep}, 30(10), 1349--1361.
\item Bonnet, M.~H., \& Arand, D.~L. (1997). Heart rate variability: sleep stage,
time of night, and arousal influences. \emph{Electroencephalography and Clinical
Neurophysiology}, 102(5), 390--396.
\item Brookshire, G., et al. (2024). Data leakage in deep learning studies of
translational EEG. \emph{Frontiers in Neuroscience}, 18, 1373515.
\item Dehghani, A., et al. (2019). A quantitative comparison of overlapping and
non-overlapping sliding windows for human activity recognition using inertial
sensors. \emph{Sensors}, 19(22), 5026.
\item Del Pup, F., et al. (2025). The role of data partitioning on the performance
of EEG-based deep learning models in supervised cross-subject analysis.
\emph{Computers in Biology and Medicine}, 196.
\item Ehrlich, F., et al. (2024). State-of-the-art sleep arousal detection
evaluated on a comprehensive clinical dataset. \emph{Scientific Reports}, 14.
\item Kattenborn, T., et al. (2022). Spatially autocorrelated training and
validation samples inflate performance assessment of convolutional neural
networks. \emph{ISPRS Open Journal of Photogrammetry and Remote Sensing}.
\item Li, A., et al. (2020). A deep learning-based algorithm for detection of
cortical arousal during sleep. \emph{Sleep}, 43(12).
\item Lotte, F., et al. (2018). A review of classification algorithms for
EEG-based brain-computer interfaces: a 10 year update. \emph{Journal of Neural
Engineering}, 15(3).
\item McAllester, D., \& Stratos, K. (2018). Formal limitations on the measurement
of mutual information. \emph{arXiv:1811.04251}.
\item Oliveira, R., et al. (2022). Split conformal prediction and non-exchangeable
data. \emph{Journal of Machine Learning Research}, 25.
\item Quinn, C.~J., Kiyavash, N., \& Coleman, T.~P. (2015). Directed information
graphs. \emph{IEEE Transactions on Information Theory}, 61(12), 6887--6909.
\item Rao, A.~A., Fredericks, M., Dresler, M., Rebollo, I., Zeitzer, J.~M.,
Schoch, S.~F., \& Coleman, T.~P. (2025). Simultaneous stomach-brain
electrophysiology reveals dynamic coupling in human sleep. \emph{bioRxiv}
2025.11.13.686572.
\item Togo, F., Cherniack, N., \& Natelson, B. (2006). Electroencephalogram
characteristics of autonomic arousals during sleep in healthy men.
\emph{Clinical Neurophysiology}, 117(12), 2597--2603.
\end{itemize}

\end{document}
